\documentclass[11pt,a4paper]{article}
\usepackage[T1]{fontenc}
\usepackage{lmodern}
\usepackage[margin=2.5cm]{geometry}
\usepackage{amsmath,booktabs,graphicx,xcolor}
\usepackage[hidelinks]{hyperref}
% Values awaiting the final runs
\newcommand{\TBD}[1]{\textcolor{red}{\textbf{[#1]}}}
\renewcommand{\thetable}{S\arabic{table}}
\renewcommand{\thefigure}{S\arabic{figure}}
\renewcommand{\thesection}{S\arabic{section}}

\title{Supplementary Material\\[4pt]
\large Inferring effective interactions in a multispecies oral biofilm model:\\ Bayesian calibration and a testable prediction}
\author{K.~Nishioka, F.~Klempt, H.~Geisler, R.~Mukherjee, N.~Heine, K.~Doll-Nikutta,\\ M.~Stiesch, S.~P.~Szafra\'{n}ski, M.~Soleimani, P.~Junker}
\date{}

\begin{document}
\maketitle

\section{Prior boxes}
\label{sec:S_boxes}
Table~\ref{tab:S_boxes} lists the uniform prior box of every interaction entry for each condition and stage.
The initial box was set empirically to contain earlier posterior estimates and encodes no sign constraint.
An entry was widened only where more than 20\% of its posterior mass lay in the outer 5\% of the range on one side, or where this fraction increased by more than 0.10 between stages with the same likelihood; the stage was then repeated.
The last column gives $\ln V=\sum_k\ln(u_k-l_k)$, the log-volume of the box, which enters $\ln\hat Z$ and is the reason why $\ln\hat Z$ is not compared across conditions.

\begin{table}[h]
\centering
\caption{Prior boxes per condition and stage for the 15 free entries; entries widened relative to stage (i) in bold. Stage (iv) uses the stage-(iii) posterior mean $\pm4$\,SD, clipped to the stage-(iii) box. $\ln V=\sum_k\ln(u_k-l_k)$ over the free entries.}
\label{tab:S_boxes}
\small\begin{tabular}{lllll}\toprule
% placeholder: run tools/make_supplementary_tables.py
\multicolumn{5}{l}{\TBD{run tools/make\_supplementary\_tables.py}}\\

\bottomrule\end{tabular}
\end{table}

\section{Runs and convergence}
\label{sec:S_runs}
Table~\ref{tab:S_runs} lists, for every condition and stage, the number of particles and mutation steps, the number of tempering stages, the acceptance rate, the moves per parameter, and the spread over the three seeds of the maximum log-likelihood and of the posterior medians (criteria (i)--(iv) of the main text).

\begin{table}[h]
\centering
\caption{Runs of the final pipeline: particles $N_p$, mutation steps $K$, tempering stages, acceptance rate, moves per free parameter, maximum log-likelihood (best seed; spread over three seeds in parentheses) and wall-clock hours per seed.}
\label{tab:S_runs}
\small\begin{tabular}{llrrlllll}\toprule
% placeholder: run tools/make_supplementary_tables.py
\multicolumn{9}{l}{\TBD{run tools/make\_supplementary\_tables.py}}\\

\bottomrule\end{tabular}
\end{table}

\section{Bayes factor for the \textit{Fn}--\textit{Pg} interaction}
\label{sec:S_bayes}
Stage~(i) was repeated with $a_{45}=0$ for all four conditions (three seeds each).
Table~\ref{tab:S_bayes} reports $\ln\hat Z$ of both models, $\ln B=\ln\hat Z_{\mathrm{full}}-\ln\hat Z_{a_{45}=0}$, and the width of the $a_{45}$ box, which sets the Occam penalty.

\begin{table}[h]
\centering
\caption{Log-evidence of the full and the reduced ($a_{45}=0$) model at stage (i), mean $\pm$ half-range over three seeds, and $\ln B=\ln\hat Z_{\mathrm{full}}-\ln\hat Z_{a_{45}=0}$. The last two columns give the $a_{45}$ box and its log-width, which enters $\ln\hat Z_{\mathrm{full}}$ as an Occam penalty.}
\label{tab:S_bayes}
\small\begin{tabular}{lccccc}\toprule
% placeholder: run tools/make_supplementary_tables.py
\multicolumn{6}{l}{\TBD{run tools/make\_supplementary\_tables.py}}\\

\bottomrule\end{tabular}
\end{table}

\paragraph{Dysbiotic Static.}
In DS, $a_{45}$ and $a_{33}$ trade off. With $a_{45}=0$ the posterior moves almost entirely to the mode $a_{33}\approx-10$ (weight \TBD{0.83--0.87}), whereas the full model places most of its mass at $a_{33}\approx+1$ (weight of the mode $a_{33}<-5$: \TBD{0.09--0.12}).
Table~\ref{tab:S_modes} gives $a_{45}$ separately for each mode.

\begin{table}[h]
\centering
\caption{DS: $a_{45}$ by $a_{33}$ mode (stage \TBD{iv}, three seeds).}
\label{tab:S_modes}
\begin{tabular}{lccc}
\toprule
Mode & Particles & $a_{45}$ median [5\%, 95\%] & max log-likelihood\\
\midrule
$a_{33}\ge-5$ & \TBD{} & \TBD{+2.4 [+0.5, +4.4]} & \TBD{}\\
$a_{33}<-5$   & \TBD{} & \TBD{+0.8 [$-$0.2, +1.4]} & \TBD{}\\
\bottomrule
\end{tabular}
\end{table}

\section{Likelihood weights}
\label{sec:S_lambda}
Stage~(i) for DH was repeated with $\lambda_{\mathrm{Pg}}=\lambda_{\mathrm{late}}=1$ (1000 particles, 100 mutation steps, three seeds; all convergence criteria met). Table~\ref{tab:S_lambda} compares it with the weighted run.

\begin{table}[h]
\centering
\caption{DH, stage (i): weighted ($\lambda_{\mathrm{Pg}}=5$, $\lambda_{\mathrm{late}}=3$) and unweighted likelihood. Ranges over three seeds.}
\label{tab:S_lambda}
\begin{tabular}{lccc}
\toprule
 & $a_{45}$ median [5\%, 95\%] & Pg D21/D15 (measured 2.74) & RMSE\\
\midrule
Weighted   & $+2.9$ to $+3.2$ [$+0.5$ to $+1.1$, $+5.2$ to $+5.5$] & 2.4--2.9 & 0.058--0.065\\
Unweighted & $+3.0$ [$-0.06$ to $+0.17$, $+5.6$] & 1.7--3.1 & 0.059--0.068\\
\bottomrule
\end{tabular}
\end{table}

\end{document}
